\documentclass[12pt]{article}
\usepackage{natbib}
\usepackage[letterpaper, margin=1in]{geometry}
\usepackage{graphicx}
\usepackage{amssymb}
\usepackage{helvet}
\renewcommand{\familydefault}{\sfdefault}
\usepackage{setspace}

\title{Can a model transcribe French speech?\\A documentation of the model}
\date{}
\author{}

\begin{document}
    \singlespacing
    \setlength{\parindent}{36pt}
    \maketitle
    \section{Background}
    To have feedback for participants in an experiment about clarity, having a
    model to be as the listener would be a logical next step. Instead of training
    a model from scratch or even taking a speech to text model and fine tuning it,
    I figured that there has to be a model that was already trained on French
    speech.

    \section{Terminology}
    \label{sec:term} I used the term fine tuned, what is it actually. Fine
    tuning a pre-existing model is just taking a model of choice and further
    training it on a selected dataset. This makes the model be more of an expert in
    a certain portion of datatypes. So for example, the model I used that would be
    utilized is the \citet{huang2022whisper-medium-french} whisper-medium model
    for ASR (automatic speech recognition) in French model which uses OpenAI's
    medium sized whisper model \citep{radford2022whisper,radford2023robust}. The
    French ASR has been trained ``on a composite dataset comprising of over 2200
    hours of French speech audio, using the train and the validation splits of Common
    Voice 11.0, Multilingual LibriSpeech, Voxpopuli, Fleurs, Multilingual TEDx,
    MediaSpeech, and African Accented French" \citep{ardila2020common, Pratap2020MLSAL, wang-etal-2021-voxpopuli, fleurs2022arxiv, salesky2021mtedx, mediaspeech2021, gigant2022african-accented-french}.

    ASR models are a type of model that does speech to text transcription. ASR
    models ``converts a speech signal to text, mapping a sequence of audio inputs
    to text outputs" \citep{hf_transformers_asr_docs}.

    \section{Methodology}
    As mentioned in section~\ref{sec:term}, I used the whisper-medium model for
    ASR in French \citep{huang2022whisper-medium-french}. This is already a
    powerful model by itself with the amount of hours it was trained on. The model
    has $769$ million parameters to tweak which will come in handy later on.

    \subsection{Data}
    Since we have data that was already collected and have the actual transcription,
    I decided to fine tune the model further with the data. To make sure that the
    model actually learns to not include any of the annotation anomolies like \{np\},
    \{cough\}, \{sp\}, etc, I made sure to remove these. I noticed that if there
    was some discrempacies with ce and le/la when comes before a vowel so I
    removed the space that would have been put when extracting from the textgrid.
    The data can be seen here: \citet{lookitsluc1-french-cleer-speech}. I split
    the data $80\%$ to $20\%$. The training split has around 4.59 thousand elements
    and the test split has around 1.15 thousand elements.

    \subsection{Fine tuned training}
    \subsubsection{Setup}
    To make the model ready to be fine-tuned, there are some elements that were
    added on top to get the metrics to a good point. Since whisper-medium-french \citep{huang2022whisper-medium-french}
    has $769$ million parameters, I did not want to tweak the parameters that gives
    the power of the whisper model to do the speech-to-text task. To only target
    the parameters that would best to tweak in my favor, I added a PEFT (Parameter-Efficient
    Fine-Tuning) via LoRA (Low-Rank Adaptation) wrapper on top of the base model (whisper-medium-french
    \citep{huang2022whisper-medium-french}).

    LoRA freezes base model’s weights and adds an adapter on top that runs along
    side model and is task specific. PEFT is a way to lower computation times and
    tunes a small subset of parameters for the task. PEFT takes additional
    parameters and changes and tweaks them. Adds small, trainable modules or
    updates a limited set of parameters instead of the full model; Learns task-specific
    patterns without affecting the core pre-trained model; Maintains efficiency
    while enabling easy customization using techniques like adapters, LoRA, and
    prompt tuning \citep{geeksforgeeks_peft}.

    To setup the LoRA and PEFT for the whisper-medium-french \citep{huang2022whisper-medium-french},
    there are some parameters I explicitly set. The first is rank (r), which is the
    amount of parameters to be trained so ''a higher rank means the model has more
    parameters to train, but it also means the model has more learning capacity,"
    \citep{hf-peft-lora-docs-v020}. In the model I have set the rank to $8$. The lora\_alpha
    is the scaling factor which I put as 16. The bias is ``whether none, all or
    only the LoRA bias parameters should be trained" \citep{hf-peft-lora-docs-v020}
    which I set to ``none''.

    \subsubsection{Worry about fine-tuning the model}

    \bibliographystyle{plainnat}
    \bibliography{custom}
\end{document}